\section{Discussion and extensions}\label{sec:discussion}

The results expose a useful division of labor. The balanced-bisection theorem
shows that a hard cut problem is already present in a restricted
two-dimensional objective. The edit theorems then provide exact, auditable
progress within the much larger encoding-tree space. Hardness motivates
certified approximation by search; local certificates state precisely what
that search has exhausted.

Grafting enlarges the neighborhood without changing the objective. One NNI is
a particular graft, while a long-range graft can relocate a subtree across
several levels in one committed step. On the five diagnosed cases,
rotation-only search had reached a one-NNI certificate while the larger graft
neighborhood still contained decreasing moves. The path formula further
indicates how to scale the operation. Cached module volumes, cuts, and affected
LCA buckets can replace full candidate reconstruction, and candidates at a
fixed tree can be scored in parallel on CPU or GPU.

The edge--LCA derivation also suggests a broader objective class. Any objective
of the form
\[
 F_\phi(T;G)=\sum_{uv\in E}w_{uv}\,
 \phi\!\left(V_{\lca_T(u,v)},d_u,d_v;G\right)
\]
admits a graft difference obtained by subtracting the old and new LCA terms.
Whether that difference has a compact path cache depends on the sufficient
statistics used by $\phi$. Dasgupta cost and Moseley--Wang revenue are natural
test cases, but a general theorem should be claimed only after the formulas and
implementations are verified for non-equivalent objectives.

The present evidence has two boundaries. The hardness theorem concerns a
balanced height-two restriction, whereas the search algorithm operates over
unrestricted binary encoding trees. The exact formulas connect these parts
through one objective but do not turn the constrained hardness result into a
hardness theorem for unrestricted tree optimization. In addition, the five
graft repairs are diagnostic cases; the frozen disjoint suite is needed to
estimate how often the larger neighborhood helps beyond those cases.

Three theoretical questions remain especially important. First, transferring
Theorem~\ref{thm:balanced-hardness} to unconstrained two-dimensional or
unrestricted tree structural entropy requires a valid balance-forcing or
normal-form construction. Second, the number and barrier profile of NNIs
needed to simulate a graft should be bounded in terms of source--target tree
distance. Third, restricted graph families may admit polynomial algorithms or
stronger optimality certificates. These questions define the next stage of
the journal analysis rather than changing the proved claims above.
